\documentclass[../../../include/open-logic-section]{subfiles}

\begin{document}

\olfileid{sth}{ordinals}{vn}	
\olsection[ساخت فون نویمان]{ساخت اعداد ترتیبی به روش فون نویمان}

\olref[sth][ordinals][iso]{thm:woalwayscomparable} اندیشه‌ای را پیش
می‌کشد. می‌توانیم اشیای معینی به نام \emph{نوع‌های ترتیب} معرفی کنیم تا
به‌جای خوش‌ترتیب‌ها به کار روند. اگر~$\ordtype{A,<}$ را برای نوع ترتیب
خوش‌ترتیب~$\tuple{A,<}$ بنویسیم، امیدواریم دو اصل زیر را به دست آوریم:
\begin{align*}
	\ordtype{A, <} = \ordtype{B, \lessdot} &
	\text{ اگر و تنها اگر } \ordeq{\tuple{A, <}}{\tuple{B, \lessdot}}\\
	\ordtype{A, <} < \ordtype{B, \lessdot}&
	\text{ اگر و تنها اگر }\ordeq{\tuple{A, <}}{\tuple{B_b, \lessdot_b}}\text{ برای یک }b \in B
\end{align*}
افزون بر این، شاید بتوانیم نوع‌های ترتیب را \emph{به‌صورت مجموعه‌هایی
معین} معرفی کنیم، درست همان‌گونه که اعداد طبیعی را به‌صورت مجموعه‌هایی
معین معرفی می‌کنیم.

رایج‌ترین راه انجام این کار---و روشی که ما در پیش خواهیم گرفت---تعریف این
نوع‌های ترتیب با استفاده از مجموعه‌های خوش‌ترتیب \emph{معیار} است. فون
نویمان نخستین بار این مجموعه‌های معیار را معرفی کرد:

\begin{defn}
مجموعهٔ~$A$ \emph{متعدی} است اگر و تنها اگر
$(\forall x \in A)x \subseteq A$. سپس~$A$ یک \emph{عدد ترتیبی} است اگر و
تنها اگر~$A$ متعدی باشد و~$\in$ آن را خوش‌ترتیب کند.
\end{defn}
\noindent
در ادامه برای اعداد ترتیبی از حروف یونانی استفاده خواهیم کرد. از تعریف
بی‌درنگ نتیجه می‌شود که اگر~$\alpha$ یک عدد ترتیبی باشد، آنگاه
$\tuple{\alpha,\in_\alpha}$ یک خوش‌ترتیب است، که در آن
$\in_\alpha = \Setabs{\tuple{x,y} \in \alpha^2}{x \in y}$. پس با اندکی
مسامحه در نمادگذاری می‌توانیم صرفاً بگوییم خود~$\alpha$ یک خوش‌ترتیب است.

نخستین چند عدد ترتیبی ما چنین‌اند:
\[
	\emptyset, \{\emptyset\},
	\{\emptyset, \{\emptyset\}\},
	\{\emptyset, \{\emptyset\}, \{\emptyset, \{\emptyset\}\}\}, \ldots
\]
توجه خواهید کرد که این‌ها همان نخستین اعداد ترتیبی‌ای هستند که در اصل
موضوع نامتناهی، یعنی در تعریف فون نویمان از~$\omega$، با آن‌ها روبه‌رو
شدیم (بنگرید به~\olref[sth][z][infinity-again]{sec}). این تصادفی نیست.
تعریف فون نویمان از اعداد ترتیبی، اعداد طبیعی را عدد ترتیبی تلقی می‌کند،
ولی اعداد ترتیبی ترامتناهی را نیز مجاز می‌شمارد.

مثل همیشه، اکنون می‌توانیم بپرسیم: آیا این‌ها \emph{واقعاً} اعداد ترتیبی
هستند؟ یا فون نویمان صرفاً مجموعه‌هایی در اختیارمان گذاشته است که
می‌توانیم آن‌ها را اعداد ترتیبی \emph{تلقی کنیم}؟ بحث‌هایی که دربارهٔ این
پرسش می‌توان داشت شبیه بحث‌هایی است که در~\olref[sfr][rel][ref]{sec}،
\olref[sfr][arith][ref]{sec}،
\olref[sfr][infinite][dedekindsproof]{sec} و~\olref[sth][z][nat]{sec}
داشتیم؛ پس بیش از این بر نکته اصرار نمی‌ورزیم. در عوض، از این پس عبارت
«اعداد ترتیبی» را صرفاً برای سخن‌گفتن از «اعداد ترتیبی فون نویمان» به کار
خواهیم برد.

\end{document}
